\newpage
\section{Arbeitstitel}

Der Titel soll die drei konstitutiven Elemente der geplanten Arbeit in ihrer 
Verbindung sichtbar machen: das zu entwickelnde Artefakt (Vorgehensmodell), 
die eingesetzte Technologie (UI-basierte Robotic Process Automation) sowie den 
spezifischen Anwendungskontext, für den das Modell ausgelegt sein soll 
(API-lose Legacy-Umgebungen). Die Bezeichnung \emph{Vorgehensmodell} markiert 
dabei den methodisch-konstruktiven Charakter der Arbeit und grenzt das Vorhaben 
von rein deskriptiven oder evaluativen Studien ab \autocite{Hevner2004}. 
Sie signalisiert, dass das Ziel nicht die Beschreibung eines vorgefundenen 
Phänomens ist, sondern die Konstruktion eines handlungsleitenden Artefakts 
im Sinne des Design Science Research \autocite{Peffers2007}.

Der Zusatz \emph{UI-basiert} soll im Titel bewusst erhalten bleiben, da er 
die technologische Schnittstelle präzisiert, über die RPA in dem untersuchten 
Kontext operiert: Software-Roboter interagieren ausschließlich über die 
grafische Benutzeroberfläche, ohne auf systemseitige Programmierschnittstellen 
zurückgreifen zu können \autocite{VanDerAalst2018}. Diese Eingrenzung ist 
nicht trivial, da der Begriff \emph{RPA} in der Literatur mitunter auch 
API-gestützte Automatisierungsformen einschließt \autocite{Wewerka2023}, 
weshalb die Spezifizierung zur Abgrenzung des Untersuchungsgegenstands 
notwendig erscheint.

Der Begriff \emph{API-lose Legacy-Umgebungen} verweist schließlich auf den 
organisatorischen und technischen Kontext, der dem Vorhaben seine praktische 
Relevanz verleiht: Unternehmen, die Systeme betreiben, die keine standardisierten 
Schnittstellen anbieten, stehen vor besonderen Integrationshürden, für die 
bestehende Implementierungsmodelle bislang keine hinreichend kontextspezifischen 
Antworten liefern \autocite{HermJaniesch2022}. % QUELLE PRÜFEN
Der Arbeitstitel beansprucht keine Endgültigkeit; er spiegelt den aktuellen 
Erkenntnisstand zu Beginn der Arbeit wider und wird im Abstimmungsprozess 
mit dem Betreuer gegebenenfalls präzisiert.

% Leitfragen (FOM-Gliederungshilfe):
% - Aussagekraeftiger Titel. Nur ein ARBEITSTITEL (nicht der finale Anmeldetitel).
% - Essenz in wenigen Worten. Passung zur Arbeit, wissenschaftlicher Anspruch und
%   Verstaendlichkeit beachten.

% TODO: Arbeitstitel formulieren.
